\subsection{Proof of the contraction theorem}
\label{app:tensor}

Throughout the proof of Theorem~\ref{thm:tensor-contraction},
$\log$ denotes a base-two logarithm and $\ln$ denotes a natural
logarithm.

\subsubsection{The coefficient ring and exact arithmetic}
\label{subsec:tensor-arithmetic}

Write $\lambda=\lceil\log_2 s\rceil$ and $e=4\lambda$.
Choose a monic polynomial $f\in\mathbb Z[U]$ of degree $e$, with
coefficients in $\{0,1\}$ below its leading coefficient, whose reduction
modulo two is irreducible.  This can be done deterministically in
polynomial time in $s$ by enumerating the $2^e$ candidates and testing
divisibility over $\mathbb F_2$ by every monic polynomial of positive
degree at most $e/2$. Finite fields admit an irreducible polynomial of
every positive degree~\cite{vzGG13}.
Put
\[
 \mathcal O=\mathbb Z_{(2)}[U]/(f),\qquad
 \mathcal K=\mathbb Q[U]/(f),\qquad
 \mathcal A=\mathcal O[V]/(V^\ell),\qquad
 \mathcal A_p=\mathcal A/(2^p).
\]
Here $\mathbb Z_{(2)}$ consists of rational numbers with odd
denominator.  Irreducibility modulo two implies irreducibility over
$\mathbb Q$, so $\mathcal K$ is a field.  For a nonzero element of
$\mathcal K$, let $v_2$ be the minimum of the ordinary two-adic
valuations of its coordinates in $1,U,\ldots,U^{e-1}$.
Reduction of $\mathcal O$ modulo two is the field $\mathbb F_{2^e}$.
Factoring out the minimum coordinate valuation therefore shows that
\[
 v_2(xy)=v_2(x)+v_2(y)\qquad(x,y\ne0).
\]
Every element of $\mathcal O$ that is nonzero modulo two is a unit because
the determinant of its multiplication matrix is invertible in
$\mathbb Z_{(2)}$, so the inverse matrix has entries in that ring.

Arithmetic in $\mathcal A_p$ takes polynomial time and space in
$s,p,\ell$.  We use the bases $1,U,\ldots,U^{e-1}$ and
$1,V,\ldots,V^{\ell-1}$, and represent each coordinate by its integer
residue between zero and $2^p-1$.  Inverting a scalar unit can be done
by inversion modulo two followed by Newton lifting, or by its
multiplication matrix, using exact polynomial and integer
operations~\cite{vzGG13}.  Divisibility and precision for
elements of $\mathcal A$ are always coordinatewise.

The binary-coordinate lifts of the $2^e$ residue-field elements give
interpolation points with pairwise unit differences.  Lagrange
interpolation at any required number of these points is therefore
valid over $\mathcal A_p$.  We will use fewer than $s^4$ points.
A contraction whose initial entries belong to $\mathcal B_{P_0,\ell}$
belongs to that subring after computation in $\mathcal A_{P_0}$,
because it is a polynomial in those entries.

When the precision increases, we lift each tensor entry separately.
For an entry originally specified modulo $2^p$, its \emph{canonical
lift} is the element of $\mathcal A$ obtained from the integer
representatives of all its coordinates. Canonical lifting need not preserve
addition or multiplication. Each recursive call contracts its specified
tensors modulo the precision of that call.  To define its deformation,
it first makes this entrywise choice of lifts.  All subsequent
evaluations in that call refer to these same lifts.
More explicitly, each tensor description is tagged with its creation
precision $p_c$ ($P_0$ for an original tensor).  At every later
precision, its entry is evaluated modulo $2^{p_c}$ and those same
integer representatives are used.  Precisions strictly increase
along a recursive path, so these representatives are also the
canonical representatives at every later precision.  In particular,
a packet tensor retains its creation-time representatives throughout the recursion.

\subsubsection{Parameters and collision packing}
\label{subsec:tensor-packing}

All estimates below hold for sufficiently large $s$, with the
threshold depending only on the fixed constants.  A fixed finite
range of smaller sizes is handled by direct summation. The parameter
inequalities can be tested directly.  Set
\[
 H=\left\lceil\frac{s}{\lfloor\sqrt\lambda\rfloor}\right\rceil,
 \qquad r=64,\qquad \varepsilon=\frac1{128\cdot1024},
 \qquad \Delta=\lceil s^{9/10}\rceil.
\]
A call with $M\le H$ edges is a leaf and is evaluated directly.
At an active call at depth $d$, with the root at depth zero, use
\begin{equation}\label{eq:tensor-parameters}
 u=u_d=\left\lfloor\frac{\varepsilon\lambda}{2^d}\right\rfloor,
 \qquad a=a_d=2\sum_{j=0}^d u_j,\qquad
 a_{-1}=0,\qquad L=2^a,\qquad D_{\rm arr}=r^{2a}.
\end{equation}
We will arrange that every child has at most $M/8$ edges and maximum
degree at most $\Delta$.  These assertions imply, on every active path,
\[
 H<M\le Cs\,8^{-d},\qquad
 d\le\frac{\log_2\lambda}{6}+O_C(1).
\]
Consequently
\begin{equation}\label{eq:tensor-parameter-bounds}
 u=\Theta_C(\lambda/2^d),\quad
 \lambda^{4/5}\le u\le M,\quad
 a\le4\varepsilon\lambda,\quad
 L,D_{\rm arr}\le s^{1/1000}
\end{equation}
for all sufficiently large $s$.  In particular, at each active call
\begin{equation}\label{eq:tensor-packing-conditions}
 M\ge2D_{\rm arr}^{\,2}\Delta,\qquad
 M\ge L D_{\rm arr},\qquad 2L\le\Delta.
\end{equation}
For example, the right side of the first inequality is
$O(s^{9/10+2/1000})=o(H)$.  The other inequalities follow in the
same way.

\begin{lemma}[Packing into packets]\label{lem:tensor-packets}
Under \eqref{eq:tensor-packing-conditions}, label the vertices on
each side of the network by elements of $[r]^a$.  There is a
deterministically constructible labeling and an embedding of the
edges into packets of $L$ slots with the following properties.
\begin{enumerate}
\item There are at most $3M/L$ packets.  Each edge occupies the same
      slot at its row and column ends. Unused slots have no ports.
\item In every packet and every coordinate, the unordered pair of
      row and column symbols is fixed.  When these symbols differ,
      bit $i$ of the slot number specifies their orientation in
      coordinate $i$, for $0\le i<a$.
\item For every coordinate and every symbol, at most $3M/r$ real
      ports on each side belong to vertices with that symbol.
\end{enumerate}
The construction uses polynomial time and space in $s+\mathcal L$.
\end{lemma}
\begin{proof}
For each pair of labels let $X_\gamma$ be the number of edges with
that pair.  Independent uniform vertex labels give
\[
 \mathbb E\sum_\gamma (X_\gamma-M/D_{\rm arr})^2\le2M\Delta.
\]
Indeed, expansion expresses this expectation as a sum over ordered
pairs of edges of their label-agreement probability minus
$1/D_{\rm arr}$.  Pairs with disjoint endpoints contribute zero.
At most $2M\Delta$ ordered pairs share an endpoint, including equal
edges, and each contributes at most one.
Conditional expectation gives a labeling with this bound.  To
implement it, try each of the $r^a$ labels of the next vertex and
compute the conditional expectation by enumerating pairs of edges.
Each probability concerns at most four vertices, so this is an
exact polynomial-time calculation.  This is the usual
conditional-expectation method~\cite{Raghavan88}.

Every load is at most $M/D_{\rm arr}+\sqrt{2M\Delta}$.
Pad all pair-label classes to a common number $H_{\rm pad}$ of
slots, the maximum load rounded up to a multiple of $L$.
The assumptions imply
\[
 H_{\rm pad}\le M/D_{\rm arr}+\sqrt{2M\Delta}+L
              \le3M/D_{\rm arr}.
\]
Fix the unordered pair of symbols in each coordinate.  If $a'$
coordinates have unequal symbols, their orientations give
$2^{a'}$ pair labels.  Partition their padded slots into packets
using $2^{a-a'}$ slots from each of these labels per packet.
This is possible because $H_{\rm pad}$ is divisible by $L$.
Number a packet's positions from $0$ to $L-1$.  Its $i$th bit
specifies the orientation when the $i$th pair is unequal and the copy
when the pair is equal.  Use one orientation convention
throughout each packet.

There are $H_{\rm pad}D_{\rm arr}/L\le3M/L$ packets.  In a fixed
coordinate, a given symbol occurs in exactly $1/r$ of all padded
pair labels on either side.  Its number of real ports is therefore
at most $H_{\rm pad}D_{\rm arr}/r\le3M/r$.
\end{proof}

\subsubsection{A Cauchy deformation and integral local coordinates}
\label{subsec:tensor-deformation}

Fix an active call, its packets, and its entrywise lifted tensors.
Use
\[
 S=\{-L,-L+1,\ldots,L-1\},\qquad
 \mathcal C(T)=\prod_{t\in S}(T-t),\qquad \alpha=2^{a+1}.
\]
Within a packet, for its row slot $x$ and column slot $y$, put
\begin{equation}\label{eq:tensor-cauchy}
 W_{xy}(T)=\frac{\mathcal C(T)}
                  {\mathcal C'(0)(T-(y-x))}.
\end{equation}
The factor in the denominator cancels a numerator factor, since
$y-x\in S$.  Thus this is a polynomial over $\mathbb Q$ of degree
$2L-1$.  Set entries between different packets to zero and restrict
to the real ports.  Under the original edge identification,
$W(0)=I$, since the diagonal entries are one and each other entry has
an uncancelled zero factor.

Let $A(Y)$ and $B(X)$ be the products of the lifted row and column
tensors, respectively, on independent port assignments
$Y,X\in\{0,1\}^M$.  Write $l_0(z)=1-z$ and $l_1(z)=z$.  Define
\begin{equation}\label{eq:tensor-deformed-contraction}
 F(T)=\sum_{X,Y\in\{0,1\}^M} A(Y)B(X)
                  \prod_{\text{row ports }x}l_{Y_x}((W(T)X)_x).
\end{equation}
Its coefficients lie in $\mathcal K[V]/(V^\ell)$, its degree is at
most $M(2L-1)$, and $F(0)$ is the contraction of the lifted tensors.
In particular, reducing $F(0)$ modulo the call's precision gives
the required answer.

\begin{lemma}\label{lem:tensor-integrality}
For every $t\in S$, the entries of $W(t+\alpha z)$ belong to
$\mathcal O[z]$, and $F(t+\alpha z)$ belongs to $\mathcal A[z]$.
Given the $e$ coordinates of $\zeta$ modulo $2^p$, each entry
$W_{xy}(t+\alpha\zeta)\bmod 2^p$ can be computed using polynomially
many bit operations in $s,p$.
\end{lemma}
\begin{proof}
The grid $S$ is a complete residue system modulo $2^{a+1}$.
For $t\ne q$ in $S$, $v_2(t-q)\le a$, and the counts of nonzero
differences at every valuation are independent of $t$.  Hence
$v_2(\mathcal C'(t))=v_2(\mathcal C'(0))$.
In
\[
 \frac{\mathcal C(t+\alpha z)}{\mathcal C'(0)}
  =\frac{\alpha z\prod_{q\in S\setminus\{t\}}
                         (t-q+\alpha z)}{\mathcal C'(0)},
\]
each factor indexed by $q\ne t$ is coefficientwise divisible by
$2^{v_2(t-q)}$.  After canceling these powers against the
denominator, the displayed polynomial is divisible by $2^{a+1}$.
To obtain \eqref{eq:tensor-cauchy}, omit just the factor indexed by
$y-x$. If $y-x=t$, this removes $\alpha z$ and leaves an integral
polynomial. Otherwise $\alpha z$ remains and the omitted factor has
two-adic content $2^{v_2(t-y+x)}$ with $v_2(t-y+x)\le a$, so the
result is still divisible by two. Substitution in
\eqref{eq:tensor-deformed-contraction} proves the assertion for $F$.

For the evaluation, first divide each factor
$t-q+\alpha z$ with $q\ne t$ by the known integer $2^{v_2(t-q)}$.
Both coefficients of the resulting linear polynomial are integers.
Cancel the remaining explicit powers of two against the denominator
as above, and invert its odd part modulo $2^p$. We can now evaluate
the product directly at the supplied residue of $\zeta$, including
$\zeta=0$. The input has $ep$ bits, and the product has $O(L)$
factors. All coordinate operations take polynomial time in $s,p$.
\end{proof}

\subsubsection{Tapered Hermite recovery}
\label{subsec:tensor-hermite}

The next lemma isolates the precision estimate.  All power-series
truncations in its statement are ordinary truncations in $z$.

\begin{lemma}[Tapered Hermite recovery]\label{lem:tapered-hermite}
Let integers $u,M$ satisfy $1\le u\le a$ and $M\ge u$, and put
$h=\lceil3M/u\rceil$.
For $t\in S\setminus\{0\}$ set
\[
 i=v_2(t),\qquad q_t=q(i)=h\max(0,u-i),\qquad D_t=L/2^i.
\]
Suppose $F\in(\mathcal K[V]/(V^\ell))[T]$ has degree at most
$M(2L-1)$, $F(0)\in\mathcal A$, and
$F(t+\alpha z)\in\mathcal A[z]$ for every $t\in S$.
For $p\ge1$, its value $F(0)$ modulo $2^p$ can be recovered from
\[
 F(t+\alpha z)\pmod{2^{p_t}},
 \qquad
 p_t=p+8hD_t+8M(a+1),\qquad q_t>0.
\]
Only the first $q_t$ coefficients of each polynomial are needed.
Recovery takes polynomial time and space in $M,L,\ell,e$ and the
precisions.
\end{lemma}
\begin{proof}
There are $L$ odd points in $S$, so
$\sum_{t\ne0}q_t\ge Lhu\ge3ML>M(2L-1)$.  Also
$h\le4M/u$, and hence $q_t\le4M$.
For each $t$ of positive multiplicity put
\[
 E_t(T)=\prod_{q\in S\setminus\{0,t\}}
                    \left(\frac{T-q}{t-q}\right)^{q_q}.
\]
Hermite interpolation gives
\begin{equation}\label{eq:tensor-hermite-origin}
 F(0)=\sum_{t:q_t>0}E_t(0)
       \left(\left[\frac{F(t+\alpha z)}
                            {E_t(t+\alpha z)}\right]_{<q_t}\right)
                            (-t/\alpha).
\end{equation}
For completeness, replace zero in the right side by an indeterminate
$T$, and $-t/\alpha$ by $(T-t)/\alpha$.  Each summand has degree
less than $\sum q_t$.  At its own node it has the prescribed jet,
and at every other node it vanishes to the required order.
The difference from $F$ is therefore divisible by
$\prod_{t\ne0}(T-t)^{q_t}$ and has smaller degree, so is zero.
This is the usual Chinese-remainder construction of Hermite
interpolation \cite[Section~5.6]{vzGG13}.

Each factor of $E_t(t+\alpha z)$ is a power of
$1+\alpha z/(t-q)$.  It is integral and has constant term one.
Its reciprocal truncated at $q_t$ is therefore integral, so
forming the quotient in \eqref{eq:tensor-hermite-origin} loses no
precision.

Fix $i=v_2(t)$ and $D=D_t$.  Directly from the product,
\begin{align}
 v_2(E_t(0))
 &=\sum_{q\ne0,t}q_q\bigl(v_2(q)-v_2(t-q)\bigr)\notag\\
 &=D\sum_{j=1}^{a-i}\frac{j}{2^j}
                           \bigl(q(i+j)-q(i)\bigr)
 \ \ge -6Dh. \label{eq:tensor-hermite-loss}
\end{align}
To justify the second equality, points of valuation less than $i$
contribute zero.  For $1\le j\le a-i$, there are $D/2^j$ points with
$v_2(q)=i+j$, and the same number with $v_2(t-q)=i+j$.
The latter have $v_2(q)=i$, and these levels exhaust the points
other than $t$ with that valuation.  These counts follow from the
complete residue grid, whose nonzero integer differences have
valuation at most $a$.  Finally,
$q(i+j)-q(i)\ge-hj$ and $\sum_{j\ge1}j^2/2^j=6$.

Evaluation of a polynomial of degree below $q_t$ at $-t/\alpha$
loses at most $4M(a+1)$ further powers of two.  Thus replacing
each quotient coefficient by any lift of its residue modulo
$2^{p_t}$ changes its contribution in
\eqref{eq:tensor-hermite-origin} by an element of $2^p\mathcal A$.
Summing these errors loses no further precision.  The reconstructed
sum is integral, since it differs from the integral value $F(0)$
by an element of $2^p\mathcal A$.

One can perform this last combination by rational arithmetic in the
specified bases.  The multipliers $E_t(0)(-t/\alpha)^j$,
$j<q_t$, are products of polynomially many rational factors of
$O(a)$ bits, counting multiplicity.  Their numerator and denominator
lengths are therefore polynomial in $M,L,a$.  There are polynomially
many such coefficients, so even a product of their denominators is
an adequate polynomial-bit common denominator.  After addition,
cancel its powers of two and reduce the resulting integral
coordinates modulo $2^p$.  This proves the algorithmic assertion.
\end{proof}

At a call with input precision $p$, apply this lemma with the
parameters in \eqref{eq:tensor-parameters}.  For every active $t$,
obtain the entire polynomial $F(t+\alpha z)$ modulo $2^{p_t}$ by
evaluation at $M(2L-1)+1$ distinct binary-coordinate lifts of
residue-field elements and Lagrange interpolation.  Lemma~\ref{lem:tensor-integrality} permits this
interpolation over $\mathcal A_{p_t}$.  The number of points is
less than $s^4$ by \eqref{eq:tensor-parameter-bounds}.
The evaluation algorithm below passes precision $p_t$ to its
children.

The input and child precisions satisfy
\begin{equation}\label{eq:tensor-precision-bound}
 p\le s\lambda^{O(1)}2^{a_{d-1}},
 \qquad
 p_t=O_C(MD_t/u).
\end{equation}
The first assertion holds at the root since $P_0\le Cs$.
Along a path, the increments in precision are
$8hD_t+8M(a+1)$, with $h\le4M/u$, $M\le Cs$, and
$D_t\le2^{a_d}$.  Summing them over $O(\log\lambda)$ depths
proves the first assertion for every call, including leaves.
For the second, $q_t>0$ implies $i<u$ and hence
\[
 D_t\ge2^{a_d-u}=2^{a_{d-1}+u}.
\]
Since $M>H$ and $u\ge\lambda^{4/5}$, this extra factor $2^u$
dominates every fixed power of $\lambda$.  Thus the first bound
on $p$, as well as $8M(a+1)$, is $O_C(MD_t/u)$.
The remaining increment $8hD_t$ has that bound directly.
All precisions remain polynomial in $s$ because
$a_d\le4\varepsilon\lambda$.

The taper accounts for the factor $1/u$ in this estimate.
For comparison, assign the same multiplicity $q_*=\Theta(M)$ to
every level below $u$, and zero to the remaining levels.  At the
last retained level $i=u-1$, the identity
\eqref{eq:tensor-hermite-loss} would give the loss
\[
 -q_*D_t\sum_{j=1}^{a-i}\frac{j}{2^j}=-\Theta(MD_t).
\]
Our multiplicities instead change by only $h=\Theta(M/u)$ per
level, giving loss $O(MD_t/u)$.  The cumulative choice of $a_d$
provides the additional factor $2^u$ needed to absorb inherited
precision and the evaluation loss.  Thus the smaller bound holds
at every retained level and every recursion depth.

\subsubsection{Modes for interactions between groups}
\label{subsec:tensor-modes}

Fix an active $t$ and an interpolation value $\zeta\in\mathcal O$
for $z$.  In this subsection $T=t+\alpha\zeta$ and all calculations
are modulo $2^{p_t}$.  Put $i=v_2(t)<u$ and $D=L/2^i$.
Group vertices on both sides by their symbol in coordinate $i$.
Entries of $W(T)$ whose two ports have the same group are called
local and are kept unchanged.

This choice of coordinate makes $W(t)$ local.  Indeed,
$W_{xy}(t)=0$ unless $y-x=t$.  Adding an integer of valuation $i$
flips bit $i$, so a surviving pair has different $i$th bits.
When the packet's two symbols in that coordinate differ, this
means that the row and column symbols agree. Equal pairs of symbols
already lie in one group.  Moving from $t$ to $t+\alpha\zeta$ introduces
the crossing terms that we now expand.

In a packet whose two symbols in coordinate $i$ differ, a row
position $x$ and a column position $y$ are in different groups
exactly when their $i$th bits are equal.  Partition columns into
tiles
\[
 y=y_0+2^{i+1}j,\qquad 0\le y_0<2^{i+1}.
\]
The column group is constant on a tile.  For a row crossing to
that tile,
\[
 v_2(t-y_0+x)\le i.
\]
Otherwise $y_0$ would agree with $x+2^i$ modulo $2^{i+1}$,
contrary to equality of their $i$th bits.  Adding $\alpha\zeta$
does not change this valuation.  A geometric expansion consequently
gives
\begin{equation}\label{eq:tensor-mode-expansion}
 W_{xy}(T)=
 \sum_{\substack{b\ge0\\b+1<p_t}}
 \frac{\mathcal C(T)}
      {\mathcal C'(0)(T-y_0+x)}
 \left(\frac{2^{i+1}}{T-y_0+x}\right)^b j^b
 \pmod {2^{p_t}}.
\end{equation}
Indeed the first fraction has valuation at least $a+1-i\ge1$,
and the ratio raised to the $b$th power has valuation at least
one before taking that power.  The geometric remainder is
divisible by $2^{p_t}$.  These expressions are computable by
factoring the known powers of two and using only unit inverses.

For each such packet, tile, and $b$ introduce a mode variable $Z$
of weight $w_Z=b+1$, to be replaced by
\[
 L_Z(X)=\sum_{\text{real columns }y\text{ in the tile}}j^bX_y.
\]
As usual $j^0=1$, including $j=0$.
Its source is the tile's column group and its target is the other
group in that packet's coordinate pair.  Formula
\eqref{eq:tensor-mode-expansion} expresses each row form as its
local part plus a linear combination of its targeted modes.
The coefficient of $Z$ is divisible by $2^{w_Z}$.
There are at most
\begin{equation}\label{eq:tensor-mode-multiplicity}
 (3M/L)2^{i+1}=6M/D
\end{equation}
variables of any one weight.

In the product of the affine factors
$l_{Y_x}((W(T)X)_x)$, every coefficient of a mode monomial
$Z^\beta$ is divisible by $2^{\sum_Zw_Z\beta_Z}$.
Consequently only the exponent tuples
\begin{equation}\label{eq:tensor-mode-budget}
 \sum_Zw_Z\beta_Z<p_t
\end{equation}
need be considered.  The number of these tuples is
\begin{equation}\label{eq:tensor-partition-count}
 \exp\bigl(O(\sqrt{Mp_t/D})\bigr)
       =\exp\bigl(O_C(M/\sqrt u)\bigr).
\end{equation}
Here is the complete counting estimate.  For any $\theta>0$,
the generating product and \eqref{eq:tensor-mode-multiplicity}
bound the number by
\[
 \exp\left(\theta p_t+\frac{6M}{D}
                   \sum_{j\ge1}-\ln(1-e^{-\theta j})\right).
\]
Expanding the logarithm and summing the geometric series gives
\[
 \sum_{j\ge1}-\ln(1-e^{-\theta j})
  =\sum_{h\ge1}\frac1{h(e^{\theta h}-1)}
  \le\frac1\theta\sum_{h\ge1}\frac1{h^2}.
\]
Take $\theta=\sqrt{M/(Dp_t)}$ and use
\eqref{eq:tensor-precision-bound}.  All mode variables and their
exponents have polynomial-size descriptions.  Tuples satisfying
\eqref{eq:tensor-mode-budget} can be enumerated with polynomial
delay and space by ordering the variables and successively choosing
each exponent subject to the remaining budget.  Every feasible partial
choice has the all-zero continuation, so the number of visited
partial choices is at most the number of complete tuples times
the number of variables, up to an additive constant.

\subsubsection{The child networks and their descriptions}
\label{subsec:tensor-children}

For one tuple $\beta$ in \eqref{eq:tensor-mode-budget} construct
a child network for each group $g\in[r]$.
Place all current vertices in group $g$, from both sides of the
current network, on one side of the child.  Retain their ports and
their chosen lifted entries, now reduced modulo $2^{p_t}$.
On the other side put one vertex for each packet, with an edge
to every retained port of that packet.  Its indices are precisely
the row bits $Y$ and column bits $X$ of group $g$ in the packet.

To define this new packet tensor, form the product of the row
factors $l_{Y_x}$ for group $g$ in the packet, with their arguments
expressed in local column bits and formal modes.  Extract the
coefficient with exponents $\beta_Z$ for all modes targeted at
$g$ in that packet.  Multiply it by
\[
 \prod_{\substack{Z\text{ in this packet}\\
                  \text{source}(Z)=g}}L_Z(X)^{\beta_Z}.
\]
Empty products equal one. If there are no target rows, the extracted
coefficient is zero whenever a target exponent is positive. If the
source tile of $Z$ has no real column, then $L_Z(X)=0$ and its
power contributes one for $\beta_Z=0$ and zero otherwise.

Every mode has one target and one source.  Coefficient extraction
therefore factors over the packet-group row products, and the
substitution for each mode depends only on bits in its source
group.  Distributing the product in
\eqref{eq:tensor-deformed-contraction}, and then summing the
independent bits within groups, proves the exact identity
\begin{equation}\label{eq:tensor-recursion-identity}
 F(t+\alpha\zeta)
   =\sum_{\beta:\,\sum w_Z\beta_Z<p_t}
                       \prod_{g=1}^r Z_{g,\beta}
                  \pmod {2^{p_t}},
\end{equation}
where $Z_{g,\beta}$ is the contraction of the child just defined.
The omitted monomials vanish by their coefficient divisibility.

By Lemma~\ref{lem:tensor-packets}, a child has at most
$6M/r=3M/32<M/8$ edges.  Old vertex degrees are unchanged, and a
new packet vertex has at most $2L\le\Delta$ ports.  Thus the
invariants used in \eqref{eq:tensor-parameter-bounds} hold
inductively.

All new tensors are given implicitly.  To evaluate a packet
tensor entry, expand its product of at most $L$ affine row
factors.  Each factor has polynomially many terms in $s$,
including its local term.  Enumerating these choices, checking
the desired exponents, and summing the resulting terms takes
\[
 \exp(O(L\log s))
\]
time and polynomial space.  The factors $L_Z(X)^{\beta_Z}$ are
computed by ordinary repeated squaring in $\mathcal A_{p_t}$.
Because $L\le s^{1/1000}$, this entry cost is
$\exp(O(s^{9/10}))$.

A description of a new tensor consists of its packet data, the
tuple $\beta$, the interpolation value, and the creation precision.
These data and the local kernel suffice to compute its entry directly.  A retained tensor
is evaluated from its original description, or from its own
creation-time description if it was created earlier in the
recursion.  Its residue at that precision is then canonically
lifted to the current precision.  The old descriptions are
immutable and shared by reference.  Thus no composition of
entry routines builds up with the depth, and all entry routines
continue to have the required time and space bounds.
Only polynomially many vertices and polynomially much new
description data are added along a path.

\subsubsection{Time and space bounds}
\label{subsec:tensor-cost}

\begin{proof}[Proof of Theorem~\ref{thm:tensor-contraction}]
Embed the original entries into $\mathcal A_{P_0}$ and perform
the recursive algorithm above.  At an active call, for every node
of positive multiplicity and every required residue-distinct
interpolation value, compute \eqref{eq:tensor-recursion-identity}
by its child contractions.  Interpolate $F(t+\alpha z)$ and apply
Lemma~\ref{lem:tapered-hermite}.  At a leaf, sum directly over all
Boolean edge assignments.  Lemma~\ref{lem:tensor-integrality},
the recursion identity, and the recovery lemma prove correctness
by induction from the leaves.

The number of interpolation nodes, values, and groups at a call
is polynomial in $s$, independently of $\mathcal L$.
By \eqref{eq:tensor-partition-count}, the number of children of
an active depth-$d$ call is at most
\[
 s^{O(1)}\exp\left(
       O_C\left(\frac{s}{\sqrt\lambda}
                    \left(\frac{\sqrt2}{8}\right)^d\right)\right).
\]
Indeed $M\le Cs\,8^{-d}$ and
$u_d=\Theta_C(\lambda/2^d)$.  Multiplying these bounds over the
$O_C(\log\lambda)$ depths, or bounding the number of nodes on
each level in this way, gives at most
\[
 \exp\left(O_C(s/\sqrt\lambda)
                   +O_C(\lambda\log\lambda)\right)
       =\exp(O_C(s/\sqrt\lambda))
\]
calls in total.  At a leaf, direct summation uses at most $2^H$
assignments and polynomially many entry evaluations per assignment.
Its cost is bounded by a polynomial in $s+\mathcal L$ times
\[
 \exp(O_C(H)+O_{C,A}(s^{9/10}))
   =\exp(O_{C,A}(s/\sqrt\lambda)).
\]
Local arithmetic, packing, descriptions, and recovery use
polynomial time per generated term, apart from the already
bounded entry computations.  The description-dependent factor
bounds the work per call. Since the branching bound is independent of
$\mathcal L$, the total work retains a fixed polynomial factor in $s+\mathcal L$.
This proves the time bound.

For space, traverse the recursion depth first, accumulating one
sum over $\beta$ and one product over groups at a time.
Store the polynomially many Taylor data needed for a call's
recovery, its current interpolation and mode indices, and its
graph and new tensor descriptions.  Their sizes are polynomial
in $s+\mathcal L$, including the precisions by
\eqref{eq:tensor-precision-bound}.  There are only
$O_C(\log\lambda)$ stack levels.  The common-denominator
calculation in Lemma~\ref{lem:tapered-hermite} also uses only
polynomial-bit integers. Each call accumulates its contraction terms
modulo its own fixed precision.
Original entry routines use polynomial space by hypothesis,
and the new routines have the same property.
This proves the space bound.  Finally extract the base
coordinates in $\mathcal A_{P_0}$, which give the answer in
$\mathcal B_{P_0,\ell}$.
\end{proof}
